\chapter{Rings, maps, ideals, and fractions}\label{chap:M07.08}

% Topic grouping: Conventions and examples
\KEntry{K07.08.ring_definition}{Associative unital rings}
\begin{definition}
A ring is an abelian group $(R,+,0)$ with an associative multiplication,
a two-sided identity $1$, and both distributive laws
$a(b+c)=ab+ac$, $(a+b)c=ac+bc$. It is commutative if $ab=ba$.
We impose $1\ne0$ only when stated. Ring homomorphisms preserve $1$.
\end{definition}
Dummit--Foote and Herstein also allow associative rings without identity;
Artin uses the word ring for a commutative unital ring; Vinberg initially
even permits nonassociative multiplication. Each argument here requires
its stated axioms, irrespective of a source's terminology. An associative
structure with the additive and distributive axioms but without an identity
will be called a nonunital ring.

Distributivity implies $a0=0a=0$, since $a0=a(0+0)=a0+a0$
and additive cancellation applies. Also
$ab+a(-b)=a(b-b)=0$, so $a(-b)=-ab$. Applying the same argument on
the other side gives $(-a)b=-ab$, and then $(-a)(-b)=ab$.

\KEntry{K07.08.zero_ring}{The zero ring}
\begin{proposition}\label{prop:zero-ring}
A unital ring has $0=1$ if and only if it has one element.
\end{proposition}
\begin{proof}
If $0=1$, every $a$ equals $a1=a0=0$. In a one-element ring the
distinguished elements $0,1$ necessarily agree.
\end{proof}
The one-element ring is allowed, and will occur as $R/R$ in
Chapter~\ref{chap:M07.09}. It is neither a domain nor a field, since those
are nontrivial by definition. Its single element is a unit, as
$0\cdot0=1$ in that ring.

\KEntry{K07.08.basic_examples}{Arithmetic, products, and functions}
The rings $\Z,\Q,\R,\C$ carry their arithmetic structures from D02.
For a family of rings $(R_i)_{i\in I}$, define operations on tuples by
$(a+b)_i=a_i+b_i$ and $(ab)_i=a_i b_i$, with coordinate identities.
Each ring axiom holds because equality of tuples is coordinatewise and
the corresponding axiom holds in each $R_i$. This proves that
$\prod_iR_i$ is a ring. In particular the function ring $R^X$ has
$(fg)(x)=f(x)g(x)$. Functions on the empty set form a one-element ring.
If $X\ne\varnothing$ and $R$ is nontrivial, constant $0$ and $1$ differ,
so $R^X$ is nontrivial.

For $n\ge1$, define operations on residue classes modulo $n$ by integer
representatives. If $a-a'$ and $b-b'$ are multiples of $n$, then
$(a+b)-(a'+b')$ and
$ab-a'b'=a(b-b')+(a-a')b'$ are too. Thus the operations are independent
of representatives. All ring identities descend from $\Z$, giving
$\Z/n\Z$; $n=1$ gives the zero ring.

\KEntry{K07.08.units}{Units and one-sided inverses}
A unit $u$ admits $v$ with $uv=vu=1$. Write $R^\times$ for the units.
\begin{proposition}\label{prop:ring-units}
An inverse of a unit is unique, and $R^\times$ is a group.
\end{proposition}
\begin{proof}
If $v,w$ are inverses of $u$, then $v=v(uw)=(vu)w=w$.
The inverse of $uv$ is $v^{-1}u^{-1}$, as multiplication on either side
shows. Closure, identity, and associativity now give the group axioms.
\end{proof}
A left and a right inverse, if both exist, coincide:
$su=1=ut$ gives $s=s(ut)=(su)t=t$. One alone is insufficient.
On a vector space with basis $e_0,e_1,\ldots$, let
$T(e_i)=e_{i+1}$, $S(e_0)=0$, and $S(e_{i+1})=e_i$.
Then $ST=\id$, whereas $TS(e_0)=0\ne e_0$, so neither endomorphism
is a unit. Units also depend on the ambient ring: $2$ is a unit in $\Q$
but not in $\Z$, whose only units are $1,-1$.

\KEntry{K07.08.domain_field}{Domains, division rings, and fields}
A division ring is a nontrivial ring in which every nonzero element is
a unit. A field is a commutative division ring. An integral domain is a
nontrivial commutative ring in which $ab=0$ implies $a=0$ or $b=0$.
In a division ring, $a\ne0$ and $ab=0$ imply
$b=a^{-1}ab=0$, so a field is a domain. Commutativity is independent:
the real quaternion ring has $ij=k=-ji$ and
\[
 (a+bi+cj+dk)^{-1}
 =\frac{a-bi-cj-dk}{a^2+b^2+c^2+d^2}
\]
for a nonzero quaternion. The product with its conjugate, on either side,
is the displayed positive real denominator, proving invertibility.
For an associative construction, represent $a+bi+cj+dk$ by
$\begin{pmatrix}a+bi&c+di\\-c+di&a-bi\end{pmatrix}$ in $M_2(\C)$.
These matrices form a real linear subspace closed under multiplication,
as multiplying the matrices for $1,i,j,k$ verifies. Their inherited
associativity and identity therefore supply the asserted ring structure.

\KEntry{K07.08.zero_divisor_cancellation}{Zero divisors and cancellation}
A zero divisor is a nonzero $a$ such that $ab=0$ or $ba=0$ for some
nonzero $b$. A unit cannot be one, since its inverse cancels either
equation. If $ax=0$ implies $x=0$, then $ab=ac$ gives
$a(b-c)=0$ and hence $b=c$. Right cancellation uses the corresponding
right-hand condition. In particular a nonzero factor cancels in a domain.
Cancelling $0$ is invalid: $0b=0c$ holds for all $b,c$.
Nonzero elements need not be either units or zero divisors: $2\in\Z$
is neither.

\KEntry{K07.08.matrix_ring}{Matrix rings}
With entrywise addition and $(AB)_{ij}=\sum_k a_{ik}b_{kj}$, the
$(i,j)$ entry of both $(AB)C$ and $A(BC)$ is
$\sum_{k,l}a_{ik}b_{kl}c_{lj}$. Finite distributivity and associativity
in $R$ justify these expansions. Entrywise expansion also verifies both
distributive laws; the diagonal matrix with diagonal entries $1$ is the
identity. Thus $M_n(R)$ is a ring.
For matrix units, direct multiplication gives
$E_{ij}E_{kl}=\delta_{jk}E_{il}$. If $R$ is nontrivial and $n\ge2$,
\[
 E_{11}E_{12}=E_{12}\ne0=E_{12}E_{11},\qquad E_{11}E_{22}=0.
\]
So the ring is noncommutative and has zero divisors, even when $R$ is
commutative. For $n=1$ it is $R$; the empty matrix ring is trivial.
These operations agree with the D05 matrix conventions.

\KEntry{K07.08.subring}{Subrings}
A unital subring contains the ambient $1$ and is closed under subtraction
and multiplication. Then $0=1-1$ belongs, negatives are differences from
$0$, and addition is subtraction of a negative; the other axioms are
inherited identities. For a nonunital subring replace the $1$ condition
by nonemptiness: a chosen $s$ gives $s-s=0$, and the same argument applies.
Thus $2\Z$ is a nonunital subring of $\Z$. It contains no identity of its
own: an identity $2k$ would satisfy $(2k)2=2$, forcing $2k=1$.

% Topic grouping: Maps and ideals
\KEntry{K07.08.ring_homomorphism}{Homomorphisms and equivalences}
A unital ring homomorphism $f:R\to S$ preserves addition, multiplication,
and $1$. Additivity gives $f(0)=f(0)+f(0)$, so $f(0)=0$;
then $f(-a)=-f(a)$ follows from $a+(-a)=0$. Repeated addition and
negation give $f(n1_R)=n1_S$ for $n\in\Z$.
Composites preserve the same operations. A bijective homomorphism is an
isomorphism; its inverse preserves operations, since applying $f$ to
$f^{-1}(xy)$ and to $f^{-1}(x)f^{-1}(y)$ gives the same element, and
$f$ is injective. The inverse preserves addition and $1$ in the same way.
The zero map between nontrivial rings is not unital. Herstein's use of
isomorphism for an injective map must be supplemented by surjectivity
when asserting that two entire rings are isomorphic.

\KEntry{K07.08.ideal_definition}{Left, right, and two-sided ideals}
A left ideal is an additive subgroup $I$ with $rI\subseteq I$ for every
$r\in R$. A right ideal has $Ir\subseteq I$, and an ideal without
qualification here is two-sided. For a commutative ring these coincide.
Absorption uses all elements of the ambient ring. A proper ideal cannot
be a unital subring.

The matrices whose only possibly nonzero column is column $j$ form a
left ideal: left multiplication changes that column without creating
others. It need not be right, since $E_{ij}$ belongs but
$E_{ij}E_{ji}=E_{ii}$ does not when $i\ne j$. A one-sided ideal cannot
silently be used as a ring-quotient kernel. In the pinned library
\texttt{Ideal R} is a left ideal; noncommutative two-sidedness requires
additional data.

\KEntry{K07.08.kernel_image}{Kernels and images}
\begin{proposition}\label{prop:ring-kernel-image}
The kernel of a ring homomorphism is a two-sided ideal, and its image is
a unital subring.
\end{proposition}
\begin{proof}
If $f(a)=f(b)=0$, then $f(a-b)=0$; also
$f(ra)=f(r)f(a)=0=f(a)f(r)=f(ar)$.
For image elements, $f(a)-f(b)=f(a-b)$ and $f(a)f(b)=f(ab)$,
and $1=f(1)$ belongs to the image.
\end{proof}
Every two-sided ideal is a kernel, by the projection constructed in
Proposition~\ref{prop:quotient-projection}; no commutativity is needed.

\KEntry{K07.08.injective_kernel}{The zero-kernel test}
\begin{proposition}\label{prop:ring-zero-kernel}
A ring homomorphism $f$ is injective if and only if $\ker f=0$.
\end{proposition}
\begin{proof}
Injectivity implies $f(a)=0=f(0)$ only at $a=0$. Conversely,
$f(a)=f(b)$ implies $f(a-b)=0$, hence $a-b=0$.
\end{proof}
This argument uses only the additive-group homomorphism.

\KEntry{K07.08.ideal_generation}{Generated ideals}
\begin{proposition}\label{prop:ideal-span}
For $S\subseteq R$, the least two-sided ideal containing $S$ is
\[
 (S)=\left\{\sum_{i=1}^m r_i s_i t_i:
 m\ge0,\ r_i,t_i\in R,\ s_i\in S\right\}.
\]
It is also the intersection of all two-sided ideals containing $S$.
For commutative $R$ the terms can be written $r_i s_i$.
\end{proposition}
\begin{proof}
The empty sum gives $0$. Concatenating sums and negating the $r_i$
give additive closure; left and right multiplication are absorbed into
$r_i$ and $t_i$. Thus the displayed set is an ideal. It contains
$s=1s1$, and every ideal containing $S$ contains its terms and sums.
Hence it is least, and has the same least-property as the indicated
intersection. Commutativity gives $r_i s_i t_i=(r_i t_i)s_i$.
\end{proof}
In particular $(\varnothing)=0$. Single products $rst$ need not be closed
under addition, so finite sums cannot be omitted. A generated left ideal
uses instead sums $r_i s_i$.

\KEntry{K07.08.ideal_sum_intersection}{The ideal lattice}
Arbitrary intersections are ideals, because subtraction and absorption
preserve membership in every member; the empty intersection is $R$.
For two ideals, $I+J=\{i+j:i\in I,j\in J\}$ is an ideal: subtraction and
either ambient multiplication preserve this form. It contains $I,J$
and lies in every ideal containing them. Intersection and sum are thus
the greatest lower and least upper bound for inclusion.
A union need not be an ideal: $(2)\cup(3)$ in $\Z$ contains $2,3$ but
not $5$. A nonempty chain union is an ideal, since any two elements
occur together in the larger of their two witnessing chain members.

\KEntry{K07.08.ideal_product}{Products of ideals}
Let $IJ$ consist of finite sums $\sum i_kj_k$ with $i_k\in I,j_k\in J$.
Addition and negatives preserve this form; left multiplication uses
$ri_k\in I$, and right multiplication uses $j_kr\in J$. Hence $IJ$ is
a two-sided ideal. Each term belongs to both $I,J$, proving
$IJ\subseteq I\cap J$.
Expanding $i(j+k)$ proves $I(J+K)\subseteq IJ+IK$.
The opposite inclusion follows from $J,K\subseteq J+K$; the right
distributive law is proved with the factors on the other side.
Both $(IJ)K$ and $I(JK)$ consist of finite sums $ijk$, giving
associativity. In $\Z$, $(2)(2)=(4)\ne(2)\cap(2)$.
Without commutativity, $IJ$ and $JI$ need not agree.

\KEntry{K07.08.principal_divisibility}{Principal ideals}
For commutative $R$, $(a)=\{ra:r\in R\}$, so
\[
 b\in(a)\ \Longleftrightarrow\ a\mid b\
 \Longleftrightarrow\ (b)\subseteq(a).
\]
Every ideal $I$ of $\Z$ is principal. If nonzero, choose its least
positive element $d$. Integer division gives $a=qd+r$,
$0\le r<d$ for $a\in I$. Since $r=a-qd\in I$, minimality forces
$r=0$, hence $I=(d)$. The zero ideal is $(0)$. Positive generators
are unique, since equality of ideals makes each divide the other.

\KEntry{K07.08.unit_in_ideal}{Ideals containing units}
\begin{proposition}\label{prop:unit-ideal}
A left ideal $I$ is $R$ if and only if $1\in I$, if and only if it
contains a unit. For commutative $R$, $(a)=R$ if and only if $a$ is a unit.
\end{proposition}
\begin{proof}
Containing $1$ gives every $r=r1$. A unit $u\in I$ gives
$u^{-1}u=1\in I$; the reverse implication uses the unit $1$.
Finally $(a)=R$ means $ra=1$, which in a commutative ring also gives
$ar=1$; the converse follows from $a^{-1}a=1$.
\end{proof}

\KEntry{K07.08.field_ideals}{Fields from ideals}
\begin{proposition}\label{prop:field-ideal-criterion}
A nontrivial commutative ring is a field if and only if its only ideals
are $0$ and $R$.
\end{proposition}
\begin{proof}
In a field an ideal containing $a\ne0$ contains a unit, hence is $R$.
Conversely $(a)\ne0$ for $a\ne0$, so $(a)=R$ and $a$ is a unit.
\end{proof}
Commutativity is essential if ideals remain two-sided. For a nonzero
ideal $I$ of $M_n(F)$ choose $A\in I$ with $a_{ij}\ne0$. Then
$E_{ki}AE_{jl}=a_{ij}E_{kl}\in I$ for arbitrary $k,l$.
The inverse scalar gives every $E_{kl}\in I$, hence $1\in I$ and
$I=M_n(F)$. For $n\ge2$ this ring still has zero divisors and is not
a division ring.

\KEntry{K07.08.field_embedding}{Maps out of fields}
A unital ring homomorphism $f:F\to R$ from a field into a nontrivial
ring is injective. Its kernel is an ideal of $F$, and cannot be $F$
since $f(1)=1\ne0$. Thus it is zero by
Proposition~\ref{prop:field-ideal-criterion}, and the zero-kernel test
applies. For a nonunital map, being nonzero excludes an entire kernel
and gives the same conclusion. The map from $F$ to the zero ring is
unital but not injective, explaining the target hypothesis.

\KEntry{K07.08.finite_domain_field}{Finite domains}
A finite integral domain is a field. If $a\ne0$, the map $x\mapsto ax$
is injective by cancellation and hence surjective on the finite set $R$.
Some $b$ therefore satisfies $ab=1$; commutativity gives $ba=1$.
Finiteness is essential, as $\Z$ shows.

% Topic grouping: Fractions
\KEntry{K07.08.fraction_field}{The fraction-field construction}
\begin{theorem}\label{thm:fraction-field-construction}
Every integral domain $R$ embeds in a field $\Frac(R)$ whose elements
have the form $a/b$, $a,b\in R$, $b\ne0$.
\end{theorem}
\begin{proof}
On such pairs set $(a,b)\sim(c,d)$ when $ad=bc$.
Reflexivity and symmetry follow from equality. If $ad=bc$ and $cf=de$,
then $adf=bcf=bde$; cancellation of nonzero $d$ gives $af=be$.
Thus this is an equivalence relation. Define
\[
 \frac ab+\frac cd=\frac{ad+bc}{bd},\qquad
 \frac ab\frac cd=\frac{ac}{bd}.
\]
Denominators stay nonzero. If $ab'=a'b$ and $cd'=c'd$, then
$acb'd'=a'c'bd$, checking multiplication. Addition is checked by
\[
 (ad+bc)b'd'=ab'dd'+bb'cd'
 =a'bdd'+bb'c'd=(a'd'+b'c')bd.
\]
So both formulas are well-defined on classes.
The negative is $(-a)/b$, zero is $0/1$, and one is $1/1$.
The two parenthesizations of a three-term sum have numerator
$adf+bcf+bde$ and denominator $bdf$; three-term products have
numerator $ace$ and denominator $bdf$.
For distributivity, both
$\frac ab(\frac cd+\frac ef)$ and
$\frac{ac}{bd}+\frac{ae}{bf}$ reduce to
$(acf+ade)/(bdf)$, cancelling a common denominator factor in the latter
expression. These computations prove associativity and distributivity.
The formulas give both identities and additive inverses, and
commutativity follows from $R$.

The map $a\mapsto a/1$ preserves the operations and $1$; equality
$a/1=c/1$ forces $a=c$, so it embeds $R$.
Also $a/b=0$ if and only if $a=0$.
For $a\ne0$, $b/a$ is an inverse of $a/b$, since $ab/(ba)=1$.
The constructed ring is nontrivial and every nonzero element is a unit.
\end{proof}
Inverting a zero divisor cannot preserve an embedding: if $bd=0$ with
$d\ne0$ and $b$ became a unit, then $d$ would become zero.
General localization accommodates such collapse, but the present
domain construction does not require it.

\KEntry{K07.08.fraction_universal}{The universal property}
\begin{theorem}\label{thm:fraction-field-universal}
An injective unital ring homomorphism $f:R\to K$ from a domain into a
field extends uniquely to a homomorphism
$\widetilde f:\Frac(R)\to K$. It is injective.
Two fraction-field models are uniquely isomorphic over $R$.
\end{theorem}
\begin{proof}
For $b\ne0$ we have $f(b)\ne0$. Set
$\widetilde f(a/b)=f(a)/f(b)$. Equality $ad=bc$ gives
$f(a)f(d)=f(b)f(c)$, proving well-definedness. The two fraction
formulas, applied in $K$, verify addition, multiplication, and $1$.
Every extension must take $(b/1)^{-1}$ to $f(b)^{-1}$, so its values
are forced. If the value at $a/b$ is zero, then $f(a)=0$, then $a=0$,
then $a/b=0$, proving injectivity. Between two fraction-field models,
construct extensions in both directions. Each composite fixes $R$,
and uniqueness makes it the identity on all fractions. This proves
existence and uniqueness of the compatible isomorphism.
\end{proof}
More generally, for a homomorphism $f:R\to S$ to a commutative ring,
an extension to $\Frac(R)$ exists if and only if every $f(b)$,
$b\ne0$, is a unit of $S$. Necessity follows because $b/1$ is
invertible. For sufficiency the same formula $f(a)f(b)^{-1}$ is
well-defined by cross multiplication with units, and the fraction
formulas prove it preserves the operations; uniqueness is unchanged.
This assertion does not claim that the extension is an embedding if
$S$ is the zero ring.
